\chapter{Прямой коннектор семейного вихря со слабым дублетом}
\label{ch:v6-spectral-transition-radial-bridge-vortex-connector-gate}

\section{Постановка после радиального запрета}

Предыдущий гейт показал, что скалярное сцепление общей амплитуды
$T|H|^2$ сохраняет $H\neq0$ в ядре вихря. Остаётся проверить более
сильную возможность: содержит ли уже построенный родитель настоящую
эквивариантную стрелку, которая переносит не только норму, но и
семейную ориентацию дефекта в слабый сектор.

Нулевая гипотеза утверждает, что такая стрелка скрыта среди уже
представленных полей или одноформ $M_{300}$. Стоп-критерий состоит из
двух частей:
\begin{enumerate}
 \item все прямые пространства интертвинеров между установленными
 семейными полями и дублетом Хиггса имеют размерность ноль;
 \item существующий смешанный модуль одноформ не содержит канонически
 выбранного ненулевого элемента.
\end{enumerate}

\section{Представления двух секторов}

Действующая группа различения равна
\begin{equation}
 G=SO(3)_{\rm fam}\times SU(2)_L\times U(1)_Y.
\end{equation}
Слабый дублет имеет метки
\begin{equation}
 H:\quad (j_{\rm fam},j_L,Y)=\left(0,\frac12,\frac12\right).
\end{equation}
Установленные компоненты вихря имеют метки
\begin{align}
 B_{\rm fam}&:(1,0,0),\nonumber\\
 Q&:(2,0,0),\nonumber\\
 T&:(3,0,0).
 \label{eq:v6-direct-connector-representations}
\end{align}
Здесь $B_{\rm fam}$ является связностью, $Q$ --- семейной пятёркой, а
$T$ --- семейной семёркой. Все они синглетны относительно
$SU(2)_L\times U(1)_Y$.

Для каждого $j=1,2,3$ машинный аудит решил линейные уравнения
эквивариантности и получил
\begin{equation}
 \dim\operatorname{Hom}_{G}
 \left((j,0,0),\left(0,\frac12,\frac12\right)\right)=0.
 \label{eq:v6-direct-connector-zero-hom}
\end{equation}
Отдельно размерность инвариантных ковекторов на каждом ненулевом
семейном спине равна нулю, а пространство интертвинеров из слабого
синглета в дублет также нулевое. Результат следует и непосредственно из
леммы Шура: исходный и конечный простые модули неэквивалентны
\cite{vortex-connector-krajewski1998,vortex-connector-paschke1997}.

Кроме того, $B_{\rm fam}$ и $H$ имеют разные пространственно-временные
типы: связность является одноформой, а Хиггс --- скаляром. Поэтому
линейное отождествление было бы запрещено уже до внутренней
классификации.

\begin{theorem}[Отсутствие прямой стрелки]
Ни семейная связность, ни поле формы $Q$, ни тетраэдрическое поле $T$ не
имеют ненулевого линейного $G$-эквивариантного отображения в слабый
дублет. Прямой коннектор среди установленных полей отсутствует.
\end{theorem}

\section{Что фактически содержится в квадрате \texorpdfstring{$M_{300}$}{M300}}

Аффинный аудит уже сузил допустимые стрелки семейной цепи до
\begin{equation}
 X=\rho V,
 \qquad
 Y=\Phi I_3.
 \label{eq:v6-existing-radial-arrows}
\end{equation}
Каждое пространство стрелок одномерно. Правое ребро пропорционально
семейной единице, поэтому смешанный член полного квадрата имеет вид
\begin{equation}
 -2\Tr(XX^\dagger)|H|^2=-2T|H|^2.
\end{equation}
Если
\begin{equation}
 XX^\dagger=T\left(\frac13I_3+Q\right),
\end{equation}
то бесследовая часть исчезает:
\begin{equation}
 \Tr Q=0.
\end{equation}
Именно поэтому квадрат видит радиальную амплитуду $T$, но не ориентацию
или форму $Q$.

Нельзя возразить, что полная матричная алгебра $M_{300}$ автоматически
содержит нужную матричную единицу. Том V уже доказал, что $M_{300}$
служит окружающей алгеброй эндоморфизмов и носителем единственного следа,
но не координатной алгеброй в определяющем представлении. Без данных
$(\mathcal A_{\rm coord},\pi,D,J,\gamma)$ она не определяет
$\Omega_D^1$ и не выделяет физические стрелки
\cite{vortex-connector-krajewski1998,vortex-connector-paschke1997}.

\begin{proposition}[Матричная единица не является физической стрелкой]
Принадлежность оператора $M_{300}(\mathbb C)$ необходима для действия на
общем носителе, но недостаточна для принадлежности представленному
комплексу одноформ. Поэтому скрытый коннектор нельзя вывести только из
размерности окружающей матричной алгебры.
\end{proposition}

\section{Смешанный модуль существует, но не выбирает секцию}

В проекте всё же имеется более тонкая зацепка. Физический модуль
заряженных одноформ имеет вид
\begin{equation}
 \mathcal Y_\rho=\mathcal E_\rho\otimes\Lambda_{\rm ch},
 \qquad
 \dim_{\mathbb C}\mathcal Y_\rho=12,
 \label{eq:v6-mixed-oneform-module}
\end{equation}
где $\mathcal E_\rho$ несёт семейную одноформу, а
$\Lambda_{\rm ch}$ содержит три хиггсовских ребра $u,d,e$. Таким
образом, семейные и слабые метки действительно встречаются в одном
тензорном модуле.

Но существование модуля не выбирает поле. Совместимые связности образуют
аффинное пространство; после удаления общего направления остаются две
вещественные относительные координаты между $u,d,e$. В более полном
моритовском углу физическое блочное ограничение оставляет не менее пяти
комплексных направлений, или четырёх после центрирования. Единственный
след задаёт их метрику, но чистая положительная норма выбирает нулевую
связность, а не ненулевой коннектор.

Следовательно, формула~\eqref{eq:v6-mixed-oneform-module} является
кандидатом для двухшагового составного механизма, но не уже выведенным
бифундаментальным полем.

\section{Разрешённый квадратичный инвариант}

Квадраты каждого целого семейного спина содержат один синглет, а
$H^\dagger H$ является слабым синглетом. Поэтому оператор
\begin{equation}
 \Tr(Q^2)H^\dagger H
\end{equation}
симметрийно допустим. Однако его коэффициент в минимальном квадрате
$M_{300}$ равен нулю. Симметрийная допустимость не заменяет происхождение
из кривизны.

Добавление нового элементарного смешанного носителя потребовало бы как
минимум представления, нетривиального одновременно в двух секторах.
Например, формальные носители с семейными спинами $1,2,3$ и слабым
дублетом имели бы соответственно комплексные размерности $6,10,14$.
Ни один из них не входит в установленный список полей; их введение было
бы расширением модели.

\begin{nogocriterion}[Запрет скрытого бифундаментала]
Нельзя выбирать произвольную матричную единицу из $M_{300}$ или добавлять
поле в представлении $(j,1/2,1/2)$ и называть его уже существующим
коннектором. Требуется либо вывести его как представленную одноформу,
либо явно объявить новую модель и заново проверить аномалии, вакуум,
нормировку и спектр.
\end{nogocriterion}

\section{Следующий различающий тест}

Прямая ветвь закрыта, но смешанный модуль
$\mathcal E_\rho\otimes\Lambda_{\rm ch}$ сохраняет один последний
внутренний маршрут. Следующий гейт
\begin{sloppypar}
\path{version6_spectral_transition_morita_two_step_connector_gate}
\end{sloppypar}
должен вычислить произведения второй степени и относительную кривизну
этого модуля. Успех требует, чтобы композиция уже существующих семейной
и хиггсовской одноформ автоматически породила ненулевой
$\Tr(Q^2)H^\dagger H$ с фиксированными знаком и нормировкой. Если
композиция блочно диагональна, исчезает в junk-идеале или сохраняет
свободное направление связности, текущая архитектура будет закрыта и для
непрямого коннектора.

\begin{protocol}
\begin{enumerate}
 \item прямые интертвинеры $B,Q,T\to H$: \textbf{нулевые};
 \item существующая семейная стрелка: \textbf{$X=\rho V$};
 \item существующая правая стрелка: \textbf{$Y=\Phi I_3$};
 \item чувствительность общего квадрата к форме $Q$:
 \textbf{отсутствует};
 \item смешанный модуль физических одноформ:
 \textbf{существует, размерность $12$};
 \item каноническая ненулевая секция смешанного модуля:
 \textbf{не выбрана};
 \item минимальная моритова неоднозначность:
 \textbf{пять комплексных направлений};
 \item скрытая физическая стрелка из одной принадлежности $M_{300}$:
 \textbf{запрещена};
 \item новый бифундаментал без расширения модели:
 \textbf{не получен};
 \item прямой коннектор вихря с Хиггсом:
 \textbf{не выведен};
 \item физическое замыкание: \textbf{не достигнуто}.
\end{enumerate}
\end{protocol}

Машинный сертификат сохранён в
\path{s2t_v6_spectral_transition_radial_bridge_vortex_connector_gate.py}
и
\path{s2t_v6_spectral_transition_radial_bridge_vortex_connector_gate_results.json}.

\begin{thebibliography}{9}
\bibitem{vortex-connector-krajewski1998}
T.~Krajewski,
\emph{Classification of Finite Spectral Triples},
Journal of Geometry and Physics 28 (1998), 1--30;
arXiv:hep-th/9701081.

\bibitem{vortex-connector-paschke1997}
M.~Paschke, A.~Sitarz,
\emph{Discrete Spectral Triples and Their Symmetries},
Journal of Mathematical Physics 39 (1998), 6191--6205;
arXiv:q-alg/9612029.
\end{thebibliography}